\chapter{Models, Meshes, and Drawing}
\label{ch:models-meshes}

\section{The runtime hierarchy}

\cnaclass{Model} owns a \cnaclass{ModelBoneCollection} and a \cnaclass{ModelMeshCollection};
each \cnaclass{ModelMesh} owns a \cnaclass{ModelMeshPartCollection} and a
\cnaclass{ModelEffectCollection}; each \cnaclass{ModelMeshPart} references a shared
\cnaclass{VertexBuffer} / \cnaclass{IndexBuffer} pair and an \cnaclass{Effect}. This is the
same shape as real XNA. The reader-side contract is deliberately separate (and is covered by
\S\ref{sec:model-loading} below and Chapter~\ref{ch:content-and-assets}); at runtime, the useful
distinction is between a normally loaded, internally consistent model and a hand-built one.
The latter can use several \texttt{CNAEXT} construction/mutation hooks, so it must preserve the
invariants that a content reader normally supplies. The exact boundary, including the remaining
intentional C++ safety deviations from FNA, matters more than a blanket claim of zero gaps.
\cnaclass{Model} exposes \texttt{CopyAbsoluteBoneTransformsTo},
\texttt{CopyBoneTransformsFrom} / \texttt{To}, and a convenience \texttt{Draw(world, view,
projection)}; one small, intentionally-kept deviation from FNA is that CNA's
\texttt{Copy*BoneTransforms*} methods loop by \texttt{Bones.Count} rather than trusting the
caller's array length the way FNA does --- strictly safer against an oversized destination
array, never a compatibility problem for correctly-sized callers.

\cnaclass{ModelBone} carries \texttt{Name}, \texttt{Index}, a get/set \texttt{Transform},
its \texttt{Parent}, and a \texttt{Children} collection built via a \texttt{CNAEXT}
\texttt{AddChild}. CNA also exposes setters for a \cnaclass{ModelMesh}'s
\texttt{BoundingSphere} and \texttt{ParentBone}; FNA reserves both setters for its internal
reader. They are useful for CNA hand construction and the XNB reader, but turn reader-populated
metadata into caller-maintained state, which is exactly where the two model loaders and the
draw contract below diverge sharply.

\section{A worked example: two ways to draw a loaded model}

The simplest path is the convenience method already named above:

\begin{lstlisting}
model.Draw(worldMatrix, camera.View(), camera.Projection());
\end{lstlisting}

This is correct for a loaded, renderable model with at least a root bone, even if it has no
interesting hierarchy to speak of. A real, multi-bone model
(a character, a vehicle with independently-animated wheels) needs each mesh drawn with its
\emph{own} absolute bone transform composed into world space first --- the classic pattern,
grounded directly in the real \texttt{CopyAbsoluteBoneTransformsTo} and
\texttt{IEffectMatrices} signatures:

\begin{lstlisting}
std::vector<Matrix> boneTransforms(model.getBonesProperty().getCountProperty());
model.CopyAbsoluteBoneTransformsTo(boneTransforms);

for (ModelMesh* mesh : model.getMeshesProperty()) {
    Matrix meshWorld = boneTransforms[mesh->getParentBoneProperty()->getIndexProperty()]
                      * worldMatrix;

    for (Effect* effect : mesh->getEffectsProperty()) {
        if (auto* matrices = dynamic_cast<IEffectMatrices*>(effect)) {
            matrices->setWorldProperty(meshWorld);
            matrices->setViewProperty(camera.View());
            matrices->setProjectionProperty(camera.Projection());
        }
    }
    mesh->Draw(); // issues actual indexed GPU draws for each valid ModelMeshPart.
}
\end{lstlisting}

\texttt{boneTransforms[mesh.getParentBoneProperty()->getIndexProperty()]} is the one line
worth pausing on: it looks up the mesh's \emph{own} parent bone's absolute transform, not bone
0 or the model's own root. This used to be exactly where the former loose-JSON route bit
hardest --- every mesh once silently got the same default bone regardless of what the source data
said, so this exact loop would have rendered every part rigidly welded together. That specific
failure mode is fixed now (Task~937, below); the worked example holds for XNB, CNJ, and direct
glTF/GLB outputs today.

\section{The draw contract: a model is more than a bag of pointers}
\label{sec:model-draw-contract}

\cnaclass{Model::Draw} is deliberately a thin convenience path, not a validator or a general
scene graph. It retains one process-wide, static vector of absolute bone matrices (the same
shape as FNA's static array), grows it only when a later model has more bones, copies the
absolute transforms into it, then visits every mesh. For every distinct effect listed by that
mesh it requires the matrix interface, writes
\texttt{boneMatrix * world}, \texttt{view}, and \texttt{projection}, and only then asks the mesh
to issue its part draws. An effect in the mesh's effect collection that does not implement that
interface is therefore a hard \texttt{runtime\_error}; it is not silently ignored or drawn with
some other matrix convention.

That compact loop establishes three input invariants which a normal content load satisfies, but
which the public/extended hand-build API does not enforce:

\begin{itemize}
  \item A renderable mesh needs an applicable bone matrix. CNA deliberately treats a null
        \texttt{ParentBone} as bone index~0, whereas FNA immediately dereferences the mesh
        parent to obtain its index and fails for null. This is a sensible C++ safety
        improvement for a one-root-bone hand-built model, and the real EasyGL/Bgfx hierarchy
        pixel tests exercise it for their root mesh. It does \emph{not} make a zero-bone model
        renderable: with a positive-effect mesh, index~0 is outside an initially empty static
        vector; after some other model has grown that shared vector, it can instead reuse that
        earlier model's stale slot. FNA's static array is likewise retained, but its usual null
        \texttt{ParentBone} path fails first rather than recovering. Give every model that can
        draw at least one bone, and give each mesh a parent explicitly when it is not that root.
  \item \texttt{CopyAbsoluteBoneTransformsTo} computes slots in collection order. A child
        therefore has to name a parent whose result was already written; it neither topologically
        sorts nor detects cycles, duplicate parents, or an index that disagrees with the vector
        position. The XNB reader's malformed-graph boundary is documented in
        Chapter~\ref{ch:content-and-assets}; the same ordering rule applies to a hand-built graph.
  \item The manual loop in the previous section intentionally dereferences the mesh parent
        pointer so that an invalid attachment is visible at the call site. Unlike
        \texttt{Model::Draw}, it has no null-to-root fallback. Well-formed
        outputs of the current loaders assign a mesh parent, but an XNB shared-resource reference
        of zero can still leave it null; a manually assembled mesh must either set one or choose
        and check its root fallback before indexing \texttt{boneTransforms}.
\end{itemize}

\cnaclass{ModelMesh::Draw} has a different, narrower safety rule: it skips a part with a null
\texttt{Effect} or a non-positive \texttt{PrimitiveCount}. FNA skips only the latter, then
dereferences a null effect, so the CNA skip is another deliberate no-undefined-behaviour
deviation. It does \emph{not} validate the remaining draw fields: a positive-count part still
passes its raw vertex buffer, index buffer, offsets, vertex count, start index, and primitive
count to \cnaclass{GraphicsDevice}. A missing buffer or invalid range consequently fails at the
graphics-device/renderer boundary rather than becoming a harmless skipped part. This is why a
model's successful construction is not proof that it is drawable.

\subsection{Keep parts and effects connected through the part setter}

The mesh's \cnaclass{ModelEffectCollection} is a derived, distinct-effect view of the effects
used by its parts; it is what \cnaclass{Model::Draw} configures before
\cnaclass{ModelMesh::Draw} applies each part's current technique. The supported way to maintain
that relationship is:

\begin{lstlisting}
part.setEffectProperty(&basicEffect); // adds basicEffect to mesh.Effects if absent
// ...
part.setEffectProperty(nullptr);      // removes it only if no other part still uses it
\end{lstlisting}

The implementation first retains an old effect while any sibling part still points to it, then
removes it only for the final user; it adds a new non-null effect only if the collection does not
already contain that pointer. Thus two parts sharing one \cnaclass{BasicEffect} produce one
matrix setup in \texttt{Model::Draw}, while each part still produces its own indexed draw. The
focused part/effect unit tests cover exactly this retain-on-shared/release-on-last-user transition
and the no-duplicate case.

\cnaclass{ModelEffectCollection::Add}/\texttt{Remove} and
\texttt{ModelMesh::getEffectsPropertyMutable()} exist only as \texttt{CNAEXT} escape hatches.
Unlike FNA's internal-only collection mutation, CNA exposes them to make its content reader
possible. They neither deduplicate \texttt{Add} nor repair a collection that a caller has
manually made inconsistent with its parts: a stale extra effect is configured but never applied
by a part, and a duplicate is configured twice. Do not populate \texttt{Effects} directly in
ordinary game code; set each part's effect and retain the buffers/effects for at least as long as
the model uses their raw pointers.

\subsection{Serial use and evidence boundary}

The static bone-matrix scratch vector has no mutex or per-model ownership. Concurrent
\texttt{Model::Draw} calls, even for different models, can resize/write it simultaneously and
are a C++ data race; concurrent mutation of bones, effects, or parts has the same ordinary
object-state problem. FNA also uses one unsynchronised static scratch array, but CNA's race is
undefined behaviour rather than merely a managed-runtime ordering hazard. Render models on the
graphics thread or otherwise serialize the entire draw/mutation sequence.

The focused CPU tests cover transform multiplication, destination-size checks, constructor
parent wiring, and part-to-effect collection bookkeeping. Renderer tests go further: the EasyGL
root/hierarchy model tests read back pixels, Bgfx has equivalent two-mesh/hierarchy tests, and
the SDL\_Renderer test proves that a fully populated positive-primitive model reaches the shared
3D draw failure boundary. No located test deliberately covers zero-bone rendering, invalid
parent order/cycles, raw out-of-range part fields, direct mutable-effect desynchronisation, or
concurrent \texttt{Draw}; treat those as construction-time invariants, not tested recovery
paths.

\section{Four model asset routes, not two}
\label{sec:model-loading}

The old two-loader and \texttt{.model.json} framing no longer describes the resolver.
CNA now has four distinct model-producing paths: the wire-compatible XNB
\texttt{ModelReader}, self-contained \texttt{.cnj}, direct \texttt{.gltf}/\texttt{.glb}, and
the older Avatar-specific \texttt{.skinnedmodel.json} reader. The first three return the
ordinary \cnaclass{Model}; the fourth returns \cnaclass{SkinnedModelEXT}, a separate runtime
type that must not be mistaken for another serialization of the same graph. For an
extensionless cache miss, \cnaclass{ContentManager} looks for \texttt{name.xnb}, then a literal
\texttt{name}, then \texttt{name.cnj}, \texttt{name.gltf}, and \texttt{name.glb}. Thus a
same-named CNJ hides a glTF/GLB asset, and an XNB sibling hides both; passing an existing literal
\texttt{name.gltf} explicitly reaches that literal before the extensionless sidecar probe. The
full resolver contract is in Chapter~\ref{ch:content-and-assets}; the resulting runtime
graphs are the important distinction here.

\begin{description}
  \item[The real \texttt{.xnb} \texttt{ModelReader}]
        (\texttt{CNA::Internal::Xnb::ModelReader}) is wire-compatible with real
        XNA/MonoGame/FNA compiled model assets. It reconstructs the full bone hierarchy, assigns
        each mesh's serialized \texttt{ParentBone}, computes its serialized
        \texttt{BoundingSphere}, and resolves shared \cnaclass{VertexBuffer}/%
        \cnaclass{IndexBuffer}/\cnaclass{Effect} resources. Its exact malformed-graph, draw-field,
        tag, and shared-resource boundary is covered in Chapter~\ref{ch:content-and-assets}.
  \item[The self-contained \texttt{.cnj} \texttt{ModelTypeReader} route]
        validates a \texttt{Model} CNJ envelope and rejects \texttt{sourceFile}; the former
        \texttt{.model.json} name survives only in historical comments and test-target names,
        not in this reader's extension list. A version-2 descriptor reconstructs its complete
        parent-before-child \texttt{bones} array and attaches every mesh to the indexed
        \texttt{parentBone}; this is how the converter preserves the selected glTF scene
        hierarchy and rigid placement. A version-1 descriptor (or one without a general
        hierarchy) retains the older compatibility shape: one root plus a synthetic named child
        bone per mesh. An optional \texttt{skeleton}/\texttt{animations} sidecar pair builds
        \cnaclass{SkinningData} on \cnaclass{Model.Tag}, and morph data lives on the corresponding
        part's \texttt{Tag}.

        This route still leaves \texttt{BoundingSphere} degenerate and
        \texttt{ModelMesh.Tag} null. It allocates a fresh vertex/index buffer and built-in stock
        effect for each described mesh part; a named custom effect instead goes through the
        manager's cached \texttt{shared\_ptr<Effect>} route. This is not XNB's serialized
        shared-resource graph, even where repeated custom-effect names happen to share a cached
        object.
  \item[The direct \texttt{.gltf}/\texttt{.glb} \texttt{ModelTypeReader} route]
        parses glTF~2.0 and its buffers without producing a CNJ or binary sidecars first. It
        reuses the offline converter's import core, including sparse-accessor-safe extraction,
        topological skin ordering, material/image extraction, and morph/animation data. Its
        visible limit is deliberate: one \texttt{Load<Model>} returns only the first
        scene-scoped mesh group, with a fixed unit scale of~1.0. Convert a multi-character or
        differently-scaled file offline when the other groups or a scale factor matter.

        Direct glTF creates an identity synthetic root followed by one \cnaclass{ModelBone} per
        reachable node of the selected scene, in parent-before-child order. A rigid mesh is
        attached to its instancing node; a skinned mesh is attached to the synthetic root so the
        skin palette does not double-apply the mesh-node transform. The separate reordered skin
        hierarchy is retained as \cnaclass{SkinningData} on \cnaclass{Model.Tag}; morph targets
        and their tracks are retained on each \cnaclass{ModelMeshPart.Tag}. Extracted images are
        decoded in memory and cached by glTF image pointer during that one import, but mesh
        buffers and effects remain per primitive.
  \item[The legacy \texttt{.skinnedmodel.json} route]
        dispatches to \texttt{SkinnedModelTypeReader} and produces
        \cnaclass{SkinnedModelEXT}, the Avatar-oriented runtime described later in
        Chapter~\ref{ch:skinning-animation}. It is neither a CNJ alias nor the historical name
        of the current Model descriptor. Code loading it opts into a different public type,
        animation representation, and draw path; its continued presence is therefore a
        compatibility route, not evidence that the generic \cnaclass{Model} reader accepts a
        fourth file spelling.
\end{description}

The evidence is much stronger than the old chapter wording suggested. The focused CNJ tests
cover a real descriptor and envelope rejection, shared animation-clip resolution has three
additional tests, and the direct runtime-glTF suite has eight unconditional end-to-end cases
(textures, reversed-order skinning/animation, morph weights and cubic splines, PBR, skinned PBR,
occlusion remapping, and punctual lights), plus a Draco case when that optional dependency is
enabled. Those tests do not make the four routes interchangeable: use XNB for compiled XNA
asset fidelity, CNJ for a controlled self-contained/imported description, direct glTF/GLB for
the first supported group of a glTF~2.0 file when the resolver priority selects it, and the
legacy route only when the application actually consumes \cnaclass{SkinnedModelEXT}.

\section{Lifetime and copies: value syntax, one mutable graph}
\label{sec:model-lifetime-copy}

\cnaclass{Model} is a copy-constructible C\texttt{++} value, but copying it is not a deep model
clone. Its compiler-generated copy operation copies the two collection vectors, \texttt{Root},
and \texttt{Tag} as raw pointers, while copying its private type-erased
\texttt{shared\_ptr<void>} ownership handle. The vectors are separate little wrappers after the
copy, but every \cnaclass{ModelBone}, \cnaclass{ModelMesh}, \cnaclass{ModelMeshPart}, buffer,
effect, and the graph below them is the same object. A bone-transform or part/effect mutation
observed through one copy is therefore observed through every copy; this is shared mutable model
state, not a way to pose or material-override one instance independently.

Each normal reader intentionally fills that last ownership handle before returning. The XNB
reader's private bundle retains unique bones, meshes, and parts plus shared vertex buffers,
index buffers, and effects. The CNJ and direct-glTF routes share a corresponding bundle which
also retains their buffers, graph objects, effects, decoded textures, skinning data, and morph
data. That is why a cached or caller-held loaded \cnaclass{Model} can safely keep the raw
pointers used by its collections: the bundle itself is shared among every copy. It also explains
why CNJ/glTF's \texttt{Model.Tag} and part Tags remain valid while such a model lives --- their
\cnaclass{SkinningData} and morph objects are in that same bundle. XNB rejects non-null serialized
Tags, so it has no analogous reader-provided Tag object to retain.

The visible cache behavior follows directly from that representation:

\begin{lstlisting}
Model first  = content.Load<Model>("ship");
Model second = content.Load<Model>("ship"); // another C++ wrapper, same bones/meshes/buffers

first.CopyBoneTransformsFrom(newPose);       // second now observes the same bone transforms
content.Unload();                            // removes manager cache entries only
// first and second still retain the reader-owned graph through their shared ownership handles
\end{lstlisting}

The real XNB test checks the decisive GPU identity part of this rule: two loads expose the same
underlying \cnaclass{VertexBuffer}, not a fresh decode/upload. The generic content-cache and
\texttt{Unload()} tests separately prove cache eviction, but no located test holds a \emph{Model}
through \texttt{Unload()}, copies one and checks a shared mutation, or exercises the final
bundle destruction. The code makes the first two outcomes clear; the missing tests remain an
evidence boundary rather than permission to assume a complete shutdown protocol.

There are two qualifications to ``shared'' worth keeping precise. Reassigning the top-level
\texttt{Tag} changes only that particular C\texttt{++} wrapper's raw pointer; it does not rewrite
another wrapper's \texttt{Tag} field, although both initially point to the same reader-owned
object. In contrast, a mesh Tag or a bone/part/effect change is made on the shared pointed-to
graph and is visible through all copies. FNA returns its cached \cnaclass{Model} \emph{class}
reference, so ordinary managed assignment aliases even the top-level wrapper; CNA's value syntax
is consequently not a deep-copy compatibility promise.

\subsection{Hand-built graphs have no automatic owner}

The public three/five-argument \cnaclass{Model} constructors merely store the caller's raw
bone and mesh addresses. Likewise, meshes store raw part and graphics-device addresses; parts
store raw buffer, index-buffer, effect, and Tag addresses; and bones store raw parent/child
addresses. They allocate none of those objects and do not install an ownership bundle. The
constructors are \texttt{CNAEXT} extensions precisely because FNA creates these graph objects
through internal readers rather than asking game code to manage them.

Consequently a hand-built model is drawable only while its complete graph outlives every model
copy that uses it; its \cnaclass{GraphicsDevice} must survive for the same draw lifetime.
Stack-allocated construction is fine for a local, synchronous test, but returning the model
after its backing bones, meshes, parts, buffers, effects, or Tag objects have died leaves
dangling pointers. An advanced CNA owner can deliberately attach one aggregate lifetime object
with the \texttt{CNAEXT}
\texttt{setOwnedResources(shared\_ptr<void>)} hook, mirroring the three readers; otherwise the
application must retain its own aggregate explicitly. The hook retains only what its supplied
shared object owns---it does not make a raw graphics-device pointer, an externally assigned Tag,
or another unowned raw address safe by itself.

Finally, \cnaclass{ContentManager::Unload()} clears its generic asset and texture lookup maps;
it does not walk a loaded model or invalidate a caller-held C\texttt{++} value. A retained
reader-created model therefore keeps its bundle after unload, but still relies on the separately
unowned graphics device being alive when it draws. Drop every model copy before tearing down that
device, and do not mistake cache eviction for a safe concurrent reload or a universal resource
destructor. Chapter~\ref{ch:content-and-assets} gives the manager-wide cache, thread, and
disposal rules; this section supplies the model-specific raw-pointer consequence.

\section{Skinning: two independent systems, not one}

CNA ships two separate, unrelated systems for animated meshes, and conflating them is an
easy mistake to make. \cnaclass{AnimationPlayer} (\texttt{CNAEXT}, mirroring Microsoft's own
unshipped XNA ``Skinned Model Sample'' reference code rather than anything in the framework
assembly itself) drives a bone-track timeline: \texttt{StartClip(clip)},
\texttt{Update(TimeSpan, relativeToCurrentTime, loop=true)}, and
\texttt{GetBoneTransforms()} / \texttt{GetWorldTransforms()} / \texttt{GetSkinTransforms()} ---
the last of which is exactly what feeds \cnaclass{SkinnedEffect::SetBoneTransforms}
(Chapter~\ref{ch:stock-effects}). Its skinning data (bone count, hierarchy, bind pose,
inverse bind pose, animation clips) is stashed on a loaded \cnaclass{Model}'s own
\texttt{Tag} --- the same convention the original Microsoft sample used, since real XNA's
\cnaclass{Model} has no dedicated skinning-data property of its own to hold it in.

\cnaclass{MorphTargetEXT} (also \texttt{CNAEXT}) is a completely independent, glTF-style
morph-target blending system: CPU-side re-blend plus a full vertex-buffer re-upload on every
call to \texttt{SetMorphWeightsEXT}, a deliberate simplicity-over-throughput tradeoff.
Because glTF's ``weights'' animation channel targets a mesh instance node rather than a bone
index, this system has no relationship to \cnaclass{AnimationPlayer}'s bone timeline at all;
it supports linear, step, and real cubic-spline (Hermite) interpolation, and attaches itself
via a loaded \cnaclass{ModelMeshPart}'s own \texttt{Tag} --- the same attachment convention,
applied to a different object in the hierarchy, for a different purpose.

A third, still more separate system exists for the Avatar subsystem specifically:
\texttt{SkinnedModelEXT} (its own \texttt{.skeleton.bin} / \texttt{.clip.bin} binary formats) is
deliberately not built on \cnaclass{Model} / \cnaclass{ModelBone} / \cnaclass{ModelMesh} at all,
has its own full test coverage, and is covered alongside the rest of the Avatar system in
this book.

\subsection{A worked example: driving AnimationPlayer, adapted from its own test suite}
\label{sec:animation-player-worked-example}

\texttt{AnimationPlayerTests.cpp} builds its fixtures around the smallest rig that can
actually exercise bone-hierarchy composition: two bones, where bone~1 is a child of bone~0
offset by $(0,1,0)$ in bind pose, and bone~0's own track moves it from the origin to
$(2,0,0)$ over one second. Adapted directly from that fixture:

\begin{lstlisting}
SkinningData data;
data.BoneCount = 2;
data.SkeletonHierarchy = {-1, 0};                              // bone 0 = root, bone 1's parent = 0
data.BindPose = {Matrix::getIdentityProperty(),
                  Matrix::CreateTranslation(Vector3(0, 1, 0))};
data.InverseBindPose = {Matrix::getIdentityProperty(), Matrix::getIdentityProperty()};

BoneTrackEXT track;
track.BoneIndex = 0;
track.Keys.push_back(Keyframe{System::TimeSpan::FromSeconds(0.0), Vector3(0, 0, 0)});
track.Keys.push_back(Keyframe{System::TimeSpan::FromSeconds(1.0), Vector3(2, 0, 0)});

AnimationClip clip;
clip.Duration = System::TimeSpan::FromSeconds(1.0);
clip.Tracks.push_back(track);
data.AnimationClips["Move"] = clip;

AnimationPlayer player(data);
player.StartClip(data.AnimationClips["Move"]);

// Every frame:
player.Update(gameTime.getElapsedGameTimeProperty(), /*relativeToCurrentTime=*/true);
skinnedEffect.SetBoneTransforms(player.GetSkinTransforms());
\end{lstlisting}

Two details this fixture makes concrete rather than abstract: \texttt{SkeletonHierarchy} is a
flat parent-index array (\texttt{-1} means ``no parent, this is a root''), the same
representation glTF and most DCC export pipelines use, not a pointer-based tree; and
\texttt{GetWorldTransforms()} versus \texttt{GetSkinTransforms()} answer two different
questions --- the former is each bone's absolute transform (what \S\ref{sec:model-loading}'s
manual mesh-draw loop above would want for a non-skinned mesh attached to a bone), the latter
is pre-multiplied by each bone's own \texttt{InverseBindPose}, which is specifically the form
\cnaclass{SkinnedEffect::SetBoneTransforms} expects, since GPU skinning needs the bone's
\emph{delta} from bind pose, not its absolute pose.

\subsection{A worked example: a two-target morph, adapted from its own test suite}

\texttt{MorphTargetEXTTests.cpp} builds every one of its cases around the same small,
hand-crafted fixture, which doubles as a clear, minimal illustration of how the whole system
fits together: a single stride-32 triangle (three vertices, each Position/Normal/TextureCoordinate)
as the base pose, and two morph targets --- target~0 pushes every vertex by $+1$ in $Z$ with no
normal change; target~1 moves only vertex~0 by $+2$ in $X$ and tilts its normal toward $+X$:

\begin{lstlisting}
MorphTargetDataEXT morph;
morph.BaseVertexBytes = BuildBaseTriangleBytes(); // 3 verts, stride-32, +Z normal, origin plane
morph.Stride = 32;

morph.PositionDeltas.push_back({Vector3(0, 0, 1), Vector3(0, 0, 1), Vector3(0, 0, 1)}); // target 0
morph.NormalDeltas.emplace_back();                                    // no normal delta, target 0

morph.PositionDeltas.push_back({Vector3(2, 0, 0), Vector3(0, 0, 0), Vector3(0, 0, 0)}); // target 1
morph.NormalDeltas.push_back({Vector3(1, 0, 0), Vector3(0, 0, 0), Vector3(0, 0, 0)});

morph.Weights = {0.0f, 0.0f}; // both targets off -- base pose

part.setTagProperty(&morph); // ModelMeshPart::Tag, the same attachment convention SkinningData
                              // uses on Model::Tag, but one level down the hierarchy
\end{lstlisting}

Driving both targets to full weight at once is the case worth reasoning through by hand, because
it is where an easy mental model (``morph targets are independent'') breaks in a way the real
math does not: \emph{additive} combination means vertex~0 (the only vertex both targets touch)
receives \textbf{both} deltas simultaneously, while vertices~1 and~2 (untouched by target~1)
receive only target~0's:

\begin{lstlisting}
SetMorphWeightsEXT(part, {1.0f, 1.0f}); // re-blends CPU-side, re-uploads the vertex buffer

// Vertex 0: base (0,0,0) + target 0's (0,0,1) + target 1's (2,0,0) = (2, 0, 1)
// Vertex 1: base (1,0,0) + target 0's (0,0,1) + target 1's zero-for-this-vertex = (1, 0, 1)
\end{lstlisting}

The blended normal at vertex~0 is the second detail worth internalizing, because it is a place
\texttt{BlendMorphTargetsEXT} does strictly more work than a naive weighted sum would: the base
normal $(0,0,1)$ plus target~1's full-weight delta $(1,0,0)$ sums to $(1,0,1)$, which has length
$\sqrt{2}$, not~1 --- \cnaclass{BlendMorphTargetsEXT} renormalizes every blended normal back to
unit length before returning, exactly because a weighted sum of unit vectors is not itself a
unit vector in general. \texttt{TextureCoordinate} (and, by the same rule, \texttt{BlendWeight}/
\texttt{BlendIndices} on a rigged mesh part) is copied from the base pose completely unchanged
by any of this --- only \texttt{Position} and \texttt{Normal} ever participate in morph blending.
Passing a weight vector whose length does not match \texttt{PositionDeltas.size()} throws rather
than silently truncating or ignoring the extra entries, verified directly by the test suite's
own \texttt{WrongWeightCountThrows} case.

Per-frame animated weights layer directly on top of this via \texttt{EvaluateMorphWeightsEXT}
and a \texttt{MorphWeightTrackEXT} attached to \texttt{MorphTargetDataEXT::WeightTrack} --- the
same three interpolation modes glTF's own ``weights'' animation channel supports (\texttt{LINEAR},
\texttt{STEP} hold-last-value, and real Hermite \texttt{CUBICSPLINE} when a keyframe's tangents
are populated), evaluated independently of \cnaclass{AnimationPlayer}'s own bone-track timeline,
consistent with this system's complete independence from bone-based skinning noted above:

\begin{lstlisting}
const auto weights = EvaluateMorphWeightsEXT(morph.WeightTrack, animationTimeSeconds);
SetMorphWeightsEXT(part, weights);
\end{lstlisting}

\section[A worked example: naming an arbitrary root bone]{A worked example: naming an
arbitrary root bone, and how the fix was proven to work}

\cnaclass{Model}'s hand-build constructor had a real, previously-undocumented gap versus FNA,
confirmed by reading FNA's own \texttt{Model.cs} directly: FNA's constructor never sets
\texttt{Root} at all --- it is real XNA's own \texttt{ModelReader} that assigns
\texttt{model.Root = bones[rootBoneIndex]} \emph{afterward}, where \texttt{rootBoneIndex} can
name any bone in the model, not necessarily the first one. CNA's original 4-argument
constructor (the one that also accepts \texttt{meshParentBones}) had no equivalent parameter at
all, silently defaulting \texttt{Root} to \texttt{bones[0]} unconditionally --- correct for the
common case, but unable to represent a hand-built model whose true root bone is not the first
entry in its own \texttt{bones} vector.

The fix is a fifth, defaulted constructor parameter, additive-only so no existing call site's
behavior changes:

\begin{lstlisting}
Model(GraphicsDevice* graphicsDevice,
      std::vector<ModelBone*> bones,
      std::vector<ModelMesh*> meshes,
      std::vector<ModelBone*> meshParentBones,
      std::size_t rootBoneIndex = 0);   // CNAEXT -- FNA's ModelReader assigns this
                                         // externally; this is the public equivalent
                                         // for a hand-built Model.
\end{lstlisting}

One edge case is worth reasoning through explicitly, because a naive implementation gets it
wrong: an empty \texttt{bones} vector must leave \texttt{Root} as \texttt{nullptr} regardless of
the requested index, rather than throwing on the harmless default value \texttt{0} --- the real
constructor body checks emptiness first, matching the existing 3-argument constructor's own
leniency, and only bounds-checks \texttt{rootBoneIndex} against a genuinely non-empty
\texttt{bones} vector:

\begin{lstlisting}
if (!bones_.bones_.empty())
{
    if (rootBoneIndex >= bones_.bones_.size())
        throw std::out_of_range("rootBoneIndex");
    root_ = bones_.bones_[rootBoneIndex];
}
// bones_ empty: root_ stays nullptr regardless of rootBoneIndex, matching the
// 3-argument constructor's own existing behavior for this case.
\end{lstlisting}

The real regression tests earned by this fix (\texttt{ModelTests.cpp}) are worth naming for the
verification technique they demonstrate, not just their existence:

\begin{lstlisting}
FiveArgConstructorDefaultRootBoneIndexMatchesFourArgBehavior
FiveArgConstructorHonorsNonZeroRootBoneIndex
FiveArgConstructorThrowsWhenRootBoneIndexOutOfRange
FiveArgConstructorEmptyBonesLeavesRootNullEvenWithDefaultIndex
\end{lstlisting}

The nonzero-index case uses a three-bone hierarchy, selects index~2, and checks both equality
with the expected bone and inequality with bone~0. A test that only checked equality could pass
even if the parameter were silently ignored and the implementation happened to default to the
requested bone by coincidence. The fix's own commit history records a real sabotage-and-revert
check applied to confirm these tests actually exercise what they claim to: with the constructor
temporarily edited to ignore \texttt{rootBoneIndex} entirely (falling back to the old, hardcoded
\texttt{bones[0]}), exactly the two tests that depend on honoring a non-default index failed as
predicted, while the other 24 \texttt{ModelTest.*} cases were unaffected --- the same
discriminating-power methodology this book applies throughout (a test suite's own
\emph{ability} to catch a real regression is itself worth verifying, not assumed from the fact
that it currently passes).

This fix closes the hand-build runtime-API gap only. Neither loose route calls the five-argument
constructor, and both reserve index~0 as their root. Direct glTF/GLB appends one
\cnaclass{ModelBone} per reachable selected-scene node beneath a synthetic identity root; a
generated version-2 CNJ serializes and reconstructs that same parent-before-child array. The
version-1 CNJ fallback instead appends one synthetic child per mesh. A glTF skin's independently
reordered palette hierarchy still lives in \cnaclass{SkinningData}, not in the model-bone index
space. Consequently neither route has a serialized arbitrary \texttt{rootBoneIndex} for this
constructor parameter to represent. The gap is closed for a hand-built model, but remains
deliberately inapplicable to the two loose content routes; the reader-specific distinctions are
tracked in \S\ref{sec:model-loading}, not hidden by the runtime fix.

\section{Collections are read views, but not uniform ones}
\label{sec:model-collections-contract}

The four model collections deliberately expose C++ iteration rather than FNA's managed
enumerators, but they do \emph{not} have one identical API.

\cnaclass{ModelBoneCollection} has checked integer lookup, name lookup,
\texttt{TryGetValue}, \texttt{Contains}, and iterators.
Its storage is populated by \texttt{Model} or a bone's \texttt{CNAEXT AddChild};
name lookup returns the first matching name, so names are useful handles, not enforced unique
identifiers.

The mesh collection has the same named-lookup family. Its integer lookup is checked, as is the
bone collection's, so an invalid index throws rather than invoking C++ undefined behaviour.
One small FNA divergence remains: FNA rejects a null \emph{or empty} mesh-name argument before
searching, while CNA's \texttt{std::string} API has no null spelling and accepts the empty
string. It can consequently return a deliberately empty-named mesh or simply return false.
Do not use an empty name as a portable not-supplied sentinel.

\cnaclass{ModelMeshPartCollection} has checked integer access, \texttt{Count}, and iterators ---
there is no part name, \texttt{Contains}, or \texttt{TryGetValue}. Its operator explicitly rejects
a negative index and one greater than or equal to \texttt{Count} with
\cnaclass{ArgumentOutOfRangeException}, matching the bone/mesh collections rather than exposing
raw \texttt{vector[]} behaviour.

\cnaclass{ModelEffectCollection} likewise has integer access, \texttt{Count},
\texttt{Contains}, and iterators, but no name lookup or \texttt{TryGetValue}. Its
integer operator has the same explicit negative/upper-bound checks. Its mutable operations are the
\texttt{CNAEXT} escape hatches described in \S\ref{sec:model-draw-contract}, not a fourth normal
collection-builder API.

All name searches are linear and return the first matching pointer; neither loader nor collection
constructor rejects duplicate names. The focused bone/mesh suites cover normal
name/index/iterator behaviour and missing names, while \texttt{ModelCollectionIndexTests.cpp}
checks negative, \texttt{Count}, and above-count indices for all four collection families plus
first/last identity. No located collection test covers the empty-name FNA difference. The safe
cross-collection rule is therefore simple: keep the model's ownership and topology coherent at
construction time and use name lookup only where that collection actually exposes it; invalid
integer indices fail loudly but are still caller errors.
